\chapter{Reprezentowalność i diagonalizacja}\label{ch:representability-diagonal}

\section{Co dokładnie będziemy udowadniać?}
W rozdziale~\ref{ch:arithmetization} przypisaliśmy
formułom i skończonym obliczeniom liczby. Sam kod jeszcze
nie wystarcza: potrzebujemy \emph{formuł arytmetycznych},
które opowiadają o tych kodach. Warto pamiętać, że
weryfikowanie danego obliczenia oraz udowodnienie w teorii
zdania o \emph{wszystkich} obliczeniach to różne zadania.

\begin{przyklad}[Różne poziomy wypowiedzi]
Obliczamy na kartce $2+3=5$, czyli w języku arytmetyki
$\overline2+\overline3=\overline5$.
Inna wypowiedź głosi, że \emph{istnieje dowód} tego zdania.
Po zakodowaniu dowodów drugą wypowiedź można zapisać
wewnątrz arytmetyki jako zdanie o pewnej liczbie $p$.
Jeszcze inna wypowiedź pyta, czy istnieje algorytm
rozstrzygający dla dowolnego zdania, czy ma dowód.
Te trzy pytania nie są tym samym.
\end{przyklad}

\section{Reprezentowalność}\label{sec:reprezentowalnosc}

\begin{definicja}[Reprezentowalność relacji]\label{res:7:1}
Zapis $\overline{\vec n}$ oznacza krotkę numerałów
$(\overline{n_1},\ldots,\overline{n_k})$; nie jest operacją
algebraiczną na krotce liczb.
Niech $T$ będzie teorią w języku $\mathcal L_{\mathrm{ar}}$ z Definicji~\ref{def:arith-language}, a $k\in\N$ ustaloną arnością. Relacja $R\subseteq\N^k$ jest \emph{reprezentowalna} w $T$, jeśli istnieje formuła $\rho(x_1,\ldots,x_k)$ bez innych zmiennych wolnych taka, że dla każdej krotki $\vec n\in\N^k$ zachodzi
\[
R(\vec n)\Longrightarrow T\vdash\rho(\overline{\vec n}),\qquad
\neg R(\vec n)\Longrightarrow T\vdash\neg\rho(\overline{\vec n}).
\]
\end{definicja}

\begin{definicja}[Reprezentowalność funkcji]\label{res:7:2}
Niech $T$ będzie teorią języka $\mathcal L_{\mathrm{ar}}$, a $k\in\N$. Funkcja częściowa $f:\N^k\rightharpoonup\N$ jest \emph{reprezentowalna} w $T$, jeśli istnieje formuła $\varphi_f(\vec x,y)$, której wszystkie zmienne wolne należą do $x_1,\ldots,x_k,y$, taka, że dla każdego standardowego wejścia $\vec n$, dla którego $f(\vec n)\downarrow=m$, zachodzi
\[
T\vdash\forall y\,\bigl(\varphi_f(\overline{\vec n},y)
\leftrightarrow y=\overline m\bigr).
\]
Na wejściach spoza dziedziny $f$ nie wymagamy dowodu nieistnienia wyniku. Dla funkcji totalnej warunek dotyczy każdego standardowego wejścia.
\end{definicja}

\begin{uwaga}[Metajęzyk i numerały pod znakami dowodliwości]
Liczba $g=\ulcorner G\urcorner$ jest w metateorii kodem napisu $G$.
Natomiast wewnątrz formuły arytmetycznej piszemy
$\operatorname{Bew}_T(\overline g)$, czyli podstawiamy \emph{term}
$\overline g=S^g(0)$. Dla czytelności w klasycznym zapisie
$\operatorname{Bew}_T(\ulcorner G\urcorner)$ numerał bywa domyślny;
w tym skrypcie w dowodach podkreślamy różnicę między liczbą
$g$ a jej numerałem $\overline g$.
\end{uwaga}

\begin{uwaga}[Obliczenie konkretne a totalność]\label{res:7:3}
Definicja~\ref{res:7:2} wymaga osobnego dowodu dla każdej ustalonej krotki numerałów; nie wynika z niej $T\vdash\forall\vec x\,\exists!y\,\varphi_f(\vec x,y)$. Pierwszy warunek wystarcza do diagonalizacji. Drugi jest istotnie
silniejszym \emph{wymogiem formalnym} i nie wolno go przyjmować
bez osobnego dowodu wewnątrz PA.
\end{uwaga}


\section{Przygotowanie do kodowania: chińskie twierdzenie o resztach}
Aby jedna liczba mogła pamiętać wiele kolejnych wartości,
potrzebujemy procedury, która równocześnie spełnia kilka
warunków na reszty z dzielenia. Przypomnimy i udowodnimy
potrzebną wersję twierdzenia, zanim użyjemy jej w kodowaniu.

\begin{definicja}[Względna pierwszość]
Liczby dodatnie $m,n$ nazywamy \emph{względnie pierwszymi},
jeśli ich jedynym wspólnym dodatnim dzielnikiem jest $1$.
Liczby $m_0,\ldots,m_s$ nazywamy \emph{parami względnie pierwszymi},
jeśli każde dwie różne są względnie pierwsze.
\end{definicja}

\begin{lemat}[Tożsamość Bézouta w potrzebnym przypadku]\label{lem:arith-bezout}
Jeśli dodatnie liczby całkowite $m,n$ są względnie pierwsze,
to istnieją $u,v\in\Z$ takie, że $um+vn=1$.
\end{lemat}
\begin{proof}
Rozważmy zbiór wszystkich dodatnich liczb postaci $am+bn$,
gdzie $a,b\in\Z$. Jest on niepusty, bo zawiera $m$.
Niech $d=um+vn$ będzie jego najmniejszym elementem.
Dzieląc $m$ przez $d$ z resztą, otrzymujemy
$m=qd+r$, gdzie $0\le r<d$. Ponieważ
$r=m-q(um+vn)=(1-qu)m-qvn$, reszta $r$ również
jest całkowitą kombinacją $m$ i $n$. Gdyby $r>0$,
sprzeciwiałaby się minimalności $d$. Zatem $r=0$
i $d\mid m$. Identyczny argument daje $d\mid n$.
Ponieważ $m,n$ są względnie pierwsze, $d=1$.
Stąd $um+vn=1$, zgodnie z tezą.
\end{proof}

\begin{twierdzenie}[Chińskie twierdzenie o resztach]\label{thm:arith-crt}
Niech $m_0,\ldots,m_s$ będą dodatnimi liczbami całkowitymi
parami względnie pierwszymi, a $c_0,\ldots,c_s\in\Z$.
Niech $M:=\prod_{i=0}^s m_i$. Istnieje dokładnie jedna
liczba $a\in\{0,\ldots,M-1\}$ taka, że
\[
a\equiv c_i\pmod{m_i}\qquad(i=0,\ldots,s).
\]
\end{twierdzenie}
\begin{proof}
Dla każdego $i$ połóżmy $M_i=M/m_i$.
Ponieważ $m_i$ jest względnie pierwsza z każdym innym czynnikiem
iloczynu, jest względnie pierwsza również z $M_i$.
Uzasadnia to wielokrotne zastosowanie Lematu~\ref{lem:arith-bezout}:
jeśli $au+bv=1$ i $ar+cw=1$, to po pomnożeniu tych
równości otrzymujemy $a(aur+cuw+brv)+bc(vw)=1$,
więc $a$ jest względnie pierwsza także z $bc$.
Powtarzamy ten argument dla kolejnych czynników $M_i$.
Z Lematu~\ref{lem:arith-bezout} istnieją zatem całkowite
$u_i,v_i$ takie, że $u_iM_i+v_im_i=1$.
Oznaczmy $a_0=\sum_{i=0}^s c_i u_i M_i$.
Ustalmy indeks $j$. Jeśli $i\ne j$, liczba $M_i$ jest
podzielna przez $m_j$, więc składnik $c_i u_i M_i$ daje
resztę zero z dzielenia przez $m_j$. Dla $i=j$ mamy
$u_j M_j\equiv1\pmod{m_j}$. Wynika stąd
$a_0\equiv c_j\pmod{m_j}$. Ponieważ $j$ był dowolny,
$a_0$ spełnia wszystkie wymagane kongruencje.

Niech $a$ będzie resztą z dzielenia $a_0$ przez dodatnie $M$.
Ma wtedy $0\le a<M$ i nadal spełnia te same kongruencje,
bo każde $m_i$ dzieli $M$. Dla jednoznaczności załóżmy,
że $a'$ jest drugim rozwiązaniem w przedziale $[0,M)$.
Wówczas każde $m_i$ dzieli $a-a'$. Dla dwóch względnie pierwszych liczb $m,n$ z
$m\mid z$ i $n\mid z$ wynika $mn\mid z$:
napiszmy $z=mk$, a z $um+vn=1$ otrzymamy
$k=u(mk)+v(nk)$ i stąd $n\mid k$. Powtarzając ten
argument dla kolejnych parami względnie pierwszych
modułów, widzimy, że $M\mid a-a'$. Mamy jednak $|a-a'|<M$, zatem $a=a'$.
\end{proof}

\begin{przyklad}[Dwie reszty w jednym numerze]
Chcemy znaleźć liczbę $a$ z resztą $2$ modulo $3$ i
resztą $3$ modulo $5$. Liczba $8$ spełnia oba warunki:
$8=2\cdot3+2=1\cdot5+3$. Ponadto każde inne
rozwiązanie różni się od $8$ o wielokrotność $15$.
Przechodząc do przedziału $[0,15)$, otrzymujemy więc
jednoznaczny wynik $8$. W Lemacie następnym będziemy
tak kodować nie dwie, lecz dowolnie wiele pozycji ciągu.
\end{przyklad}

\begin{lemat}[Kodowanie ciągów funkcją $\beta$]\label{res:7:4}
Funkcja $\beta(a,b,i)=a\bmod(1+(i+1)b)$ jest pierwotnie rekurencyjna. Dla każdego skończonego ciągu $c_0,\ldots,c_s$ istnieją $a,b\in\N$ takie, że $\beta(a,b,i)=c_i$ dla $i\le s$. Jej graf w modelu standardowym opisuje formuła $\Delta_0$
\[
B(a,b,i,r):=r<1+(i+1)b\ \land\
\exists q\le a\,[a=q(1+(i+1)b)+r].
\]
Jeśli wejściem jest kod $w=\langle c_0,\ldots,c_s\rangle$ niepustego ciągu
w numeracji z Definicji~\ref{res:6:1}, można wybrać parę $(a,b)$
funkcją pierwotnie rekurencyjną argumentu $w$.
\end{lemat}
\begin{proof}
Pierwotna rekurencyjność $\beta$ wynika z Lematu~\ref{res:5:8}. Wybierzmy $b=(s+1)!(1+\max_{i\le s}c_i)$ i oznaczmy $m_i=1+(i+1)b$. Mamy $c_i<m_i$. Moduły $m_0,\ldots,m_s$ są parami względnie pierwsze. Gdyby dwa z nich miały wspólny dzielnik większy od $1$, najmniejszy dzielnik większy od $1$ byłby liczbą pierwszą $p$ dzielącą oba. Dla $i<j$ otrzymalibyśmy $p\mid(m_j-m_i)=(j-i)b$. Ponieważ $m_i=1+(i+1)b$, liczba $p$ nie dzieli $b$. Pierwszość $p$ wymusza zatem $p\mid(j-i)$; ponieważ $1\le j-i\le s$, liczba $j-i$ dzieli $(s+1)!$, a więc dzieli $b$. To sprzeczność.

Z chińskiego twierdzenia o resztach istnieje $a$ spełniające $a\equiv c_i\pmod{m_i}$ dla każdego $i\le s$. Możemy wybrać $0\le a<M$, gdzie $M=\prod_{i\le s}m_i$; wtedy $\beta(a,b,i)=c_i$. Przypomnijmy, że konstrukcja reszt jest efektywna: dla $M_i=M/m_i$ wyznaczamy współczynniki Bézouta $u_iM_i+v_im_i=1$ i bierzemy resztę liczby $\sum_i c_i u_iM_i$ modulo $M$.

Wybór kodu jest nawet p.r. Z wejściowego kodu $w$
odczytujemy pierwotnie rekurencyjnie jego długość $s+1$ oraz
składniki $c_0,\ldots,c_s$ (Lemat~\ref{res:6:5}).
Obliczenie $b$ i $M$ wymaga tylko skończonych iloczynów, silni
i maksimum; są to operacje p.r.
Dla każdego $a<M$ sprawdzenie wszystkich kongruencji jest p.r.,
a na mocy udowodnionego istnienia i jednoznaczności z
Twierdzenia~\ref{thm:arith-crt} ograniczona minimalizacja
wybiera jedno $a$. Wartość na niepoprawnych kodach $w$
ustalamy jako $(0,0)$, zachowując totalność procedury.
Formuła $B$ jest $\Delta_0$, bo dzielnik jest dodatni, a iloraz $q$
nie przekracza $a$.
\end{proof}

\begin{lemat}[Arytmetyczny opis skończonego obliczenia]\label{res:7:5}
Dla ustalonego programu rejestrowego $e$ z Definicji~\ref{res:6:2} istnieje formuła $\Delta_0$ $C_e(\vec x,y,t,\vec a,\vec b)$ taka, że
\[
e(\vec n)\downarrow=m\quad\Longleftrightarrow\quad
\mathfrak N\models\exists t\,\exists\vec a\,\exists\vec b\,
C_e(\overline{\vec n},\overline m,t,\vec a,\vec b).
\]
Ponadto PA dowodzi jednoznaczności wyniku: dwa ślady o tym samym wejściu, spełniające $C_e$, mają ten sam wynik.
\end{lemat}
\begin{proof}
Konfiguracja programu składa się ze stanu i zawartości $d$ rejestrów, czyli $d+1$ liczb. Ustalamy kolejność współrzędnych $u=(u_0,u_1,\ldots,u_d)$: $u_0$ jest numerem stanu, natomiast $u_{h+1}$ jest zawartością rejestru $r_h$, $0\le h<d$. Każdą składową śladu kodujemy osobną parą $(a_h,b_h)$ z Lematu~\ref{res:7:4}, gdzie $0\le h\le d$. Przyjmijmy skróty
$B_i(u):=\bigwedge_{h=0}^d B(a_h,b_h,i,u_h)$ i
$A:=1+\sum_{h=0}^d a_h$; zapis $\exists u<A$ oznacza ograniczenie wszystkich współrzędnych krotki $u$ przez $A$.
Niech $I_e(\vec x,u)$ oznacza poprawną konfigurację początkową, $H_e(u)$ --- stan zatrzymania, a $S_e(u,v)$ --- wykonanie jednej instrukcji. Są to formuły bez kwantyfikatorów. Dla każdej z ustalonych skończenie wielu instrukcji wzrost rejestru zapisujemy równością $v_{h+1}=S(u_{h+1})$, a zmniejszenie dodatniej zawartości równością $u_{h+1}=S(v_{h+1})$; przypadek zerowy rozdzielamy warunkiem $u_{h+1}=0$. Pozostałe rejestry zachowują wartości, a numer stanu zmienia się zgodnie z instrukcją. Alternatywa tych przypadków definiuje $S_e(u,v)$ i wyraża deterministyczny pojedynczy krok bez użycia odejmowania. Formuła $I_e$ ustala stan początkowy i rejestry wejściowe, a $H_e$ rozpoznaje końcowy numer stanu; $u_1$ jest zawartością rejestru wynikowego $r_0$.

Definiujemy
\begin{align*}
C_e(\vec x,y,t,\vec a,\vec b):={}&
\exists u<A\,[B_0(u)\land I_e(\vec x,u)]\\
&\land\forall i<t\ \exists u<A\ \exists v<A\,
[B_i(u)\land B_{S(i)}(v)\land\neg H_e(u)\land S_e(u,v)]\\
&\land\exists u<A\,[B_t(u)\land H_e(u)\land u_1=y].
\end{align*}
Jest to formuła $\Delta_0$. Dla rzeczywistego śladu wszystkie odczyty $\beta(a_h,b_h,i)$ są nie większe od $a_h<A$, więc kodowanie z poprzedniego lematu dostarcza świadków. Odwrotnie, w $\N$ jednoznaczność reszty sprawia, że formuła opisuje kolejne konfiguracje tego samego obliczenia. Warunek $\neg H_e(u)$ dla $i<t$ oznacza, że ślad kończy się przy pierwszym zatrzymaniu.

Uzasadnijmy ostatnie stwierdzenie \emph{wewnątrz PA}. Arytmetyka ta dowodzi istnienia i jednoznaczności reszty przy dzieleniu przez dowolną dodatnią liczbę: istnienie otrzymujemy przez indukcję po dzielnej, a jednoznaczność z własności porządku i mnożenia. Stąd PA dowodzi, że każdy kod $(a_h,b_h)$ wyznacza dokładnie jedną wartość składowej w dowolnej chwili. Dla dwóch śladów o tym samym wejściu indukcja po wspólnym odcinku ich czasów wykazuje równość kolejnych konfiguracji: baza wynika z $I_e$, a krok z deterministyczności $S_e$. Gdyby jeden ślad kończył się wcześniej, drugi wymagałby w tej samej chwili $\neg H_e$, sprzecznie z $H_e$ w pierwszym. Czasy zatrzymania, a więc i wyniki, są równe. Argument wykorzystuje indukcję w PA; nie przenosimy go bez zmian do $Q$.
\end{proof}

\begin{twierdzenie}[Reprezentowalność funkcji rekurencyjnych]\label{res:7:6}
Każda funkcja częściowo rekurencyjna jest reprezentowalna w $Q$ w sensie Definicji~\ref{res:7:2}; dotyczy to w szczególności funkcji totalnych rekurencyjnych. Każda relacja rekurencyjna jest reprezentowalna w $Q$ w sensie Definicji~\ref{res:7:1}. Formuły te reprezentują odpowiednie funkcje i relacje także w każdym rozszerzeniu $T\supseteq Q$.
\end{twierdzenie}
\begin{proof}
Na mocy Lematu~\ref{res:6:3} wybieramy program $e$ obliczający $f$. Nie możemy powtórzyć rozumowania z Lematu~\ref{res:7:5}, ponieważ $Q$ nie ma indukcji i nie dowodzi ogólnej jednoznaczności dzielenia z resztą. Wzmocnimy więc formułę kodującą ślad o warunki potrzebne do prześledzenia \emph{każdego konkretnego} obliczenia.

Niech
\begin{align*}
D(t)&:=\forall u\le t\ \forall j<u\,(j<t),\\
U(\vec a,\vec b,t)&:=\forall i\le t\ \bigwedge_{h=0}^d\forall v\,\forall w\,
\bigl((B(a_h,b_h,i,v)\land B(a_h,b_h,i,w))\to v=w\bigr).
\end{align*}
Pierwszy warunek jest dodatkową hipotezą o parametrze czasu $t$:
jeżeli $u\le t$ i $j<u$, możemy wprost wywnioskować $j<t$,
bez przyjmowania ogólnego twierdzenia o przechodniości $<$ w $Q$.
Drugi gwarantuje jednoznaczność odczytów również wtedy,
gdy parametry kodu w pewnym modelu $Q$ są niestandardowe. Kładziemy
\[
\varphi_f(\vec x,y):=\exists t\,\exists\vec a\,\exists\vec b\,
\bigl(C_e(\vec x,y,t,\vec a,\vec b)\land D(t)\land U(\vec a,\vec b,t)\bigr).
\]
Nie zakładamy, że formuła $\varphi_f$ jest $\Sigma_1$; wymaganie to nie występuje w Definicji~\ref{res:7:2}.

Ustalmy $\vec n$ takie, że $f(\vec n)\downarrow=m$, i niech $s$ będzie rzeczywistą liczbą kroków do zatrzymania. Wybieramy standardowe kody tego śladu. Zamknięte instancje $C_e$ i $D(\overline s)$ są prawdziwymi zdaniami $\Delta_0$, zatem $Q$ je dowodzi na mocy Lematu~\ref{res:4:7}. Dla wybranych numerałów $\overline{\vec a},\overline{\vec b},\overline s$
teoria $Q$ dowodzi również $U$. Ustalmy składową $h$ i
standardową chwilę $j\le s$. Z Lematu~\ref{res:4:4} możemy
rozwinąć $i\le\overline s$ na skończoną listę wartości
$i=\overline0,\ldots,\overline s$. Po podstawieniu
$i=\overline j$ każdy warunek
$B(\overline{a_h},\overline{b_h},\overline j,v)$ implikuje
$v<\overline{M_{h,j}}$, gdzie $M_{h,j}=1+(j+1)b_h>0$.
Ten sam lemat sprowadza $v$ do jednego ze standardowych
numerałów $\overline0,\ldots,\overline{M_{h,j}-1}$.
Analogicznie świadek ilorazu spełnia $q\le\overline{a_h}$,
więc ma jedną ze skończenie wielu wartości standardowych.
Dla każdej pary $(q,v)$ rozstrzygamy w $Q$ instancję
$\overline{a_h}=q\overline{M_{h,j}}+v$ przez rachunek na
numerałach. Dla ustalonego $a_h$ i dodatniego $M_{h,j}$
dokładnie jedna para spośród par z $0\le v<M_{h,j}$
spełnia tę równość. Po złożeniu skończenie wielu dowodów
$Q$ dowodzi zatem, że dowolne dwa odczyty $v,w$ są równe.
Tak otrzymujemy zamkniętą instancję $U$ bez korzystania
z indukcji wewnątrz $Q$. Ponieważ rzeczywisty ślad spełnia
$C_e$, uzyskujemy $Q\vdash\varphi_f(\overline{\vec n},\overline m)$.

Pozostało wykazać, że \emph{dowolny} wynik wskazany przez
$\varphi_f(\overline{\vec n},y)$ jest równy $\overline m$,
choć świadkowie $t,\vec a,\vec b$ mogą być niestandardowi.
Niech $u^{(0)},\ldots,u^{(s)}$ będą konkretnymi standardowymi
konfiguracjami rzeczywistego obliczenia, zatrzymanego po
$s$ krokach. W $Q$ rozpatrujemy dowolnych świadków
$C_e\land D\land U$ i stosujemy lemat~\ref{res:4:4}
do rozłącznych przypadków $t<\overline s$,
$t=\overline s$ i $\overline s<t$.

Wyjaśnijmy wspólny mechanizm: dla każdej \emph{ustalonej}
standardowej chwili $j$, do której obowiązują warunki śladu,
można w $Q$ otrzymać $B_j(\overline{u^{(j)}})$. W chwili $0$
wynika to z $I_e(\overline{\vec n},u)$ i $B_0(u)$
w $C_e$: początkowa konfiguracja jest wymuszona przez
równości definiujące $I_e$. Jeśli mamy odczyt rzeczywistej
konfiguracji w chwili $j$ oraz warunek przejścia $\overline j<t$,
$U$ utożsamia współrzędne użyte w obu odczytach kodu
$B_j$. Trzeba jeszcze uzasadnić zastosowanie $U$ do
następnej chwili. Dla ustalonego standardowego $j$ teoria $Q$
dowodzi $(\overline j<t)\to(\overline{j+1}\le t)$:
z (Q8) otrzymujemy $t=z+S(\overline j)$, a (Q3)
daje $z=0$ albo $z=S(w)$. W pierwszym przypadku
rachunek na ustalonym numerale prowadzi do
$t=\overline{j+1}$; w drugim do
$t=S^{j+2}(w)=w+S(\overline{j+1})$, a więc
$\overline{j+1}<t$ na mocy (Q8). Zatem $U$ obejmuje
także odczyt $B_{j+1}$. Ponieważ $u^{(j)}$ jest ustaloną
standardową konfiguracją, po podstawieniu jej numerałów
skończona alternatywa instrukcji $S_e$ redukuje się w $Q$
do jednej właściwej gałęzi: z (Q1)--(Q7) rozstrzygamy
przypadek zerowy lub dodatni i obliczamy wszystkie
nowe współrzędne. Otrzymujemy dokładnie
$u^{(j+1)}$. Powtarzamy tę \emph{skończoną} liczbę kroków
w metateorii; budujemy odpowiedni formalny dowód w $Q$
dla każdego potrzebnego standardowego $j$, nie stosując
indukcji po dowolnym elemencie modelu $Q$.

Jeżeli $t<\overline s$, Lemat~\ref{res:4:4} daje skończoną
alternatywę $t=\overline0\lor\cdots\lor t=\overline{s-1}$.
Dla każdego jej składnika rozwijamy dokładnie tyle kroków
opisanej konstrukcji, ile wymaga dany numerał. W rozpatrywanej gałęzi $t=\overline j$ otrzymujemy
$B_j(\overline{u^{(j)}})$, a rzeczywista konfiguracja w tej chwili
nie jest końcowa, bo $s$ jest pierwszym czasem zatrzymania.
Końcowy warunek $H_e$ z $C_e$ daje więc sprzeczność.

Jeżeli $t=\overline s$, ten sam skończony argument odtwarza
konfigurację $\overline{u^{(s)}}$ w chwili zatrzymania. Ponieważ
$u_1=y$ w ostatnim warunku $C_e$, otrzymujemy
$y=\overline m$.

Wreszcie jeśli $\overline s<t$, warunek $D(t)$ pozwala
uzyskać $\overline j<t$ dla każdego standardowego $j<s$:
bierzemy $u=\overline s\le t$ i korzystamy ze standardowego
porównania $\overline j<\overline s$. Sama nierówność
$\overline s<t$ pozwala ponadto zastosować warunek
przejścia również w chwili $i=\overline s$. Skończenie
wiele wcześniejszych kroków daje odczyt konfiguracji
$\overline{u^{(s)}}$, już końcowej, natomiast przejście w tej chwili
wymaga $\neg H_e(\overline{u^{(s)}})$. Sprzeczność. Gdy $s=0$,
nie ma kroków pośrednich i wystarczy konfiguracja początkowa.

Otrzymaliśmy formalnie w $Q$ implikację
$\varphi_f(\overline{\vec n},y)\to y=\overline m$.
Ponieważ dowiedliśmy też $\varphi_f(\overline{\vec n},\overline m)$,
podstawianie równych termów daje implikację odwrotną
$y=\overline m\to\varphi_f(\overline{\vec n},y)$.
Generalizując $y$, otrzymujemy
$Q\vdash\forall y\,(\varphi_f(\overline{\vec n},y)\leftrightarrow y=\overline m)$.

Dla relacji rekurencyjnej $R$ bierzemy jej totalną funkcję
charakterystyczną i kładziemy
$\rho_R(\vec x):=\varphi_{\chi_R}(\vec x,\overline1)$.
Jeśli $R(\vec n)$, otrzymujemy dowód
$\rho_R(\overline{\vec n})$; w przeciwnym razie
reprezentowalność wartości $0$ daje dowód jej negacji,
ponieważ $Q\vdash\overline0\ne\overline1$.
Dowody z $Q$ są też dowodami w każdym rozszerzeniu $Q$.
\end{proof}

\begin{przyklad}[Dlaczego samo istnienie śladu nie wystarcza do silnej reprezentacji]
Załóżmy, że program na wejściu $2$ wykonuje dwa kroki i zwraca $5$.
Sama formuła „istnieje ślad kończący się wynikiem $y$” pokazuje,
że $Q$ dowodzi jej instancji dla $y=\overline5$.
Definicja~\ref{res:7:2} wymaga jednak więcej: $Q$ musi jeszcze
dowiedzieć, że \emph{każdy} wynik opisany przez tę formułę
jest równy $\overline5$, nawet jeśli wartości zmiennych kodujących
ślad są niestandardowe w pewnym modelu $Q$.
Z tego właśnie powodu w Twierdzeniu~\ref{res:7:6}
dodaliśmy $U$ (jednoznaczność odczytu kodu) i $D$
(legalność odwołań do wcześniejszych chwil). Rozumowanie
rozdziela trzy przypadki czasu $t$: mniejszy od rzeczywistego
czasu $2$, równy mu i większy. W pierwszym przypadku nie może
jeszcze wystąpić zatrzymanie; w drugim wynik jest równy $5$;
w trzecim kod żąda kolejnego kroku po zatrzymaniu, co daje
sprzeczność. Nie zakładamy indukcji po zmiennej $t$ w $Q$.
\end{przyklad}

\begin{lemat}[Standardowy arytmetyczny predykat dowodu]\label{res:7:7}
Dla efektywnie aksjomatyzowalnej teorii $T$ można wybrać kodowanie pełnych dowodów tak, aby formuła $\operatorname{Prf}_T(p,x)$ była klasy $\Delta_0$ i poprawnie rozstrzygała relację dowodu w modelu standardowym. Dla każdej pary numerałów $\overline p,\overline x$ teoria $Q$ dowodzi odpowiedniej instancji $\operatorname{Prf}_T$ albo jej negacji. Wówczas $\operatorname{Bew}_T(x):=\exists p\,\operatorname{Prf}_T(p,x)$ jest formułą $\Sigma_1$.
\end{lemat}
\begin{proof}
Z Lematu~\ref{res:6:5} mamy pierwotnie rekurencyjny weryfikator $V_T(d,x)$ pełnego dowodu $d$, zawierającego certyfikaty pochodzenia aksjomatów. Ustalamy program obliczający $V_T$ i dołączamy do $d,x$ kod całego śladu jego obliczenia, kończącego się wynikiem $1$.

Dane $(d,t,\vec a,\vec b)$, przy ustalonych długościach krotek zależnych od wybranego programu weryfikującego,
pakujemy iterowanym parowaniem
$\langle u,v\rangle=(u+v)^2+u$. Jest ono p.r. i różnowartościowe:
przy ustalonej sumie $u+v=s$ przyjmuje wartości w przedziale
$[s^2,s^2+s]$, rozłącznym z przedziałami dla innych sum.
Każda składowa jest nie większa od pełnego kodu $p$,
gdyż $\langle u,v\rangle\ge u,v$ dla $u,v\in\N$;
własność zachowuje iteracja. Dla ustalonej liczby składowych
wielokrotnie złożone parowanie jest określonym termem
języka $\mathcal L_{\mathrm{ar}}$: każdą parę zastępujemy
wyrażeniem $(u+v)(u+v)+u$. Niech
$\operatorname{Pack}(p;d,t,\vec a,\vec b)$
oznacza równość $p$ z tym termem, a więc formułę
bez kwantyfikatorów. Ograniczenie we wszystkich następujących
kwantyfikatorach rozumiemy współrzędnie i niezależnie:
każda współrzędna $d,t,\vec a,\vec b$ jest nie większa od $p$. Definiujemy
\[
\operatorname{Prf}_T(p,x):=\exists d,t,\vec a,\vec b\le p\,
\bigl(\operatorname{Pack}(p;d,t,\vec a,\vec b)\land
C_{V_T}(d,x,\overline1,t,\vec a,\vec b)\bigr).
\]
Ponieważ $\operatorname{Pack}$ nie ma kwantyfikatorów, $C_{V_T}$
jest $\Delta_0$, a wszystkie dodatkowe kwantyfikatory są
ograniczone przez term $p$, cała formuła jest $\Delta_0$.
Ustalone urządzenie $V_T$ jest totalnym weryfikatorem z Lematu~\ref{res:6:5}:
na wejściu $(d,x)$ zwraca $1$ dokładnie wtedy, gdy $d$
jest pełnym, poprawnym dowodem formuły o kodzie $x$.
Każdy taki dowód daje więc zakończone obliczenie weryfikatora
i kod jego śladu; z kolei ślad zakończony wynikiem $1$
może istnieć jedynie dla zaakceptowanej pary $(d,x)$.
Powstaje równoważność w \emph{modelu standardowym}; nie
przyjmujemy z niej bez dowodu ogólnej własności w modelach $Q$.
Otrzymany $p$ koduje zarówno dowód, jak i świadectwo
obliczenia, dlatego relacja $\operatorname{Prf}_T$ ma tę samą
projekcję egzystencjalną na kod zdania co
$\operatorname{Prov}_T$ z Lematu~\ref{res:6:5}.
Dla dwóch konkretnych numerałów prawdziwość lub fałszywość
zamkniętej formuły $\Delta_0$ jest dowodliwa w $Q$ na mocy
Lematu~\ref{res:4:7}. Egzystencjalne domknięcie po $p$
jest formułą $\Sigma_1$.
\end{proof}

% Schemat diagonalizacji i lemat o punkcie stałym.

\section{Schemat operacji diagonalnej}
W poniższym rysunku każda strzałka oznacza
\emph{ściśle określoną operację na skończonym napisie albo jego kodzie}.
Prawdziwość formuły nie jest tutaj jeszcze rozpatrywana.

\begin{center}
\begin{tikzpicture}[node distance=5mm,>=Stealth,
  every node/.style={font=\small},box/.style={draw,rounded corners,align=center,inner sep=5pt}]
\node[box] (a) {$\theta(x)$\\formuła};
\node[box,right=of a] (b) {$m=\ulcorner\theta\urcorner$\\kod formuły};
\node[box,right=of b] (c) {$\theta(\overline m)$\\podstawienie};
\node[box,right=of c] (d) {$d(m)$\\kod wyniku};
\draw[->] (a)--(b);
\draw[->] (b)--(c);
\draw[->] (c)--(d);
\end{tikzpicture}
\end{center}

Dla przykładu, jeśli $\theta(x)$ oznacza $x=x$, to kod $m$
oznacza napis $x=x$, a $d(m)$ oznacza kod zdania
$\overline m=\overline m$. W następnym lemacie
przeniesiemy taki zabieg z poziomu rachunku na napisach
do poziomu dowodliwych zdań arytmetycznych.

\section{Lemat o punkcie stałym (diagonalizacja)}\label{sec:diagonalizacja}
\begin{lemat}[Lemat diagonalny / o punkcie stałym]\label{res:8:1}
Niech $\psi(x)$ będzie formułą języka $\mathcal L_{\mathrm{ar}}$,
której jedyną możliwą zmienną wolną jest $x$.
Istnieje zdanie $G$ takie, że dla liczby $g:=\ulcorner G\urcorner$
zachodzi $Q\vdash G\leftrightarrow\psi(\overline g)$.
Ta sama równoważność jest dowodliwa w każdej teorii
$T\supseteq Q$ tego języka, niezależnie od efektywności jej aksjomatów. Zdanie $G$ „mówi o sobie”, że ma własność $\psi$; nie jest to dosłowna równość napisów, lecz dowodliwa równoważność.
\end{lemat}
\begin{proof}
Ustalamy wyróżnioną zmienną $x$ i definiujemy funkcję diagonalizacji
\[
d(n):=\ulcorner\theta(\overline n)\urcorner,
\]
gdy $n$ koduje formułę $\theta(x)$ bez innych zmiennych wolnych; na pozostałych wejściach kładziemy $d(n)=0$. Jest to funkcja p.r. na mocy Lematu~\ref{res:6:5}(c): podstawiamy numerał $\overline n$ w formule o kodzie $n$. Twierdzenie~\ref{res:7:6} daje formułę $\varphi_d(x,y)$ z
\[
Q\vdash\forall y\,(\varphi_d(\overline n,y)\leftrightarrow y=\overline{d(n)})
\]
dla każdego standardowego $n$.

Najpierw przemianowujemy zmienne związane w $\psi(x)$
i $\varphi_d$, jeśli mogłyby kolidować ze zmienną $y$;
następnie zastępujemy wolne wystąpienia $x$ w $\psi$ przez $y$.
Dzięki tej kolejności żadne podstawienie nie przechwytuje zmiennej.
Definiujemy $\delta(x):=\exists y\,(\varphi_d(x,y)\land\psi(y))$,
$m:=\ulcorner\delta(x)\urcorner$ i $G:=\delta(\overline m)$.
Formuła $\delta$ ma co najwyżej zmienną wolną $x$, toteż
$G$ jest zdaniem. Oznaczmy $g:=\ulcorner G\urcorner$.
Z definicji $d(m)=\ulcorner\delta(\overline m)\urcorner=g$,
więc reprezentowalność daje
\[
Q\vdash\forall y\,(\varphi_d(\overline m,y)\leftrightarrow y=\overline g).
\]
Korzystając z równoważności dla $\varphi_d(\overline m,y)$
i z reguł logiki pierwszego rzędu z równością, otrzymujemy
\[
\begin{aligned}
Q&\vdash G\leftrightarrow\exists y\,(y=\overline g\land\psi(y)),\\
Q&\vdash \exists y\,(y=\overline g\land\psi(y))
\leftrightarrow\psi(\overline g).
\end{aligned}
\]
Z przechodniości równoważności otrzymujemy
$Q\vdash G\leftrightarrow\psi(\overline g)$.
Pierwsza równoważność wynika przez zastąpienie formuły
$\varphi_d(\overline m,y)$ równoważną formułą $y=\overline g$
pod kwantyfikatorem, druga przez eliminację kwantyfikatora
egzystencjalnego i zastępowanie równych.
To jest szukana równoważność. Nie użyto żadnej hipotezy o aksjomatach $T$.
\end{proof}



\section{Dlaczego lemat diagonalny nie jest zwykłym podstawieniem?}
\begin{przyklad}[Instrukcja „podstaw własny numer”]
Rozważmy formułę $\theta(x)$ z jedną zmienną wolną.
Jej kod $m=\ulcorner\theta\urcorner$ jest liczbą naturalną.
Operacja $d(m)$ tworzy kod napisu $\theta(\overline m)$.
Dostajemy więc zależność między \emph{dwiema liczbami}:
kodem formuły i kodem jej instancji.
Lemat~\ref{res:8:1} pozwala opisać tę zależność
\emph{formułą arytmetyczną}, a następnie użyć jej
wewnątrz teorii $Q$.
Nie stwierdzamy, że zdanie jest dosłownie swoim
kodem; stwierdzamy, że pewna równoważność jest
\emph{formalnie dowodliwa}.
\end{przyklad}
\begin{uwaga}[Granica tej iteracji]
W kolejnych rozdziałach wykorzystamy zbudowaną relację
dowodu oraz lemat diagonalny do pierwszego twierdzenia
o niezupełności. Będziemy musieli oddzielnie rozważyć
niesprzeczność, $\omega$-niesprzeczność i wzmocnienie Rossera;
bez tych warunków nie wolno formułować jednej
bezwarunkowej wersji twierdzenia.
\end{uwaga}